\documentclass[UTF8,fontset=none,aspectratio=169,10pt]{beamer}
\usepackage{ads-beamer-theme}
\title{高级分布式系统}
\subtitle{第十七讲：P2P、覆盖网络、自稳定与系统综合}
\author{陈鹏飞\quad 计算机学院}
\institute{中山大学}
\date{2025--2026 学年}

\begin{document}
\begin{frame}[plain]
\begin{tikzpicture}[remember picture,overlay]
  \fill[ADSBlue] (current page.north west) rectangle ([xshift=2.4cm]current page.south west);
  \fill[ADSOrange] ([xshift=2.4cm]current page.north west) rectangle ([xshift=2.55cm]current page.south west);
  \node[anchor=west,align=left] at ([xshift=3.3cm,yshift=1.0cm]current page.west) {\fontsize{32}{38}\selectfont\color{ADSBlue}\bfseries 高级分布式系统};
  \node[anchor=west,align=left] at ([xshift=3.3cm,yshift=-0.4cm]current page.west) {\fontsize{21}{27}\selectfont\color{ADSDarkBlue}第十七讲\quad P2P 与系统综合};
  \node[anchor=west,align=left] at ([xshift=3.3cm,yshift=-2.2cm]current page.west) {\large\color{ADSGray}当节点会加入、离开、失忆甚至乱掉，系统如何继续找到数据？};
  \node[anchor=east] at ([xshift=-0.8cm,yshift=-3.2cm]current page.east) {\small\color{ADSGray}陈鹏飞\quad 中山大学\quad 2025--2026};
  \draw[ADSBlue,line width=0.8pt] ([xshift=3.3cm,yshift=-1.35cm]current page.west) -- ([xshift=12.6cm,yshift=-1.35cm]current page.west);
\end{tikzpicture}
\end{frame}

\sectionframe{一、P2P 的问题意识}{没有中央目录，也没有固定成员；位置、负载和信任都在变化。}

\begin{frame}{本讲学习目标}
\begin{enumerate}
  \item 理解 P2P 与客户端—服务器模型的差异；
  \item 区分结构化和非结构化覆盖网络，掌握 DHT 的数据定位思路；
  \item 能解释 Chord 的一致性哈希、finger table 和 churn 维护；
  \item 认识 CAN、Tapestry 的不同几何/前缀路由取舍；
  \item 用自稳定的 convergence、closure 分析系统恢复；
  \item 把 overlay、复制、自稳定连接到 AI 数据与模型基础设施。
\end{enumerate}
\vfill
\callout{主线}{把“找数据”变成“在逻辑图上路由”；把“系统乱了”变成“从任意状态回到合法状态”。}
\end{frame}

\begin{frame}{P2P 与客户端—服务器}
\begin{columns}[T,onlytextwidth]
\column{.48\linewidth}
\blocktitle{客户端—服务器}
\begin{itemize}
  \item 目录、调度和权限集中；
  \item 控制简单，中心是瓶颈和故障点；
  \item 位置通常由 DNS 或注册表解析。
\end{itemize}
\column{.48\linewidth}
\blocktitle{P2P}
\begin{itemize}
  \item 节点角色对称，资源由参与者共同提供；
  \item 自组织、分布式控制、可承受 churn；
  \item 搜索的是对象，而不是某台固定主机。
\end{itemize}
\end{columns}
\end{frame}

\begin{frame}{P2P 设计目标}
\begin{tabularx}{\linewidth}{>{\bfseries}p{3cm}XX}
\toprule
自组织 & 节点可以动态加入和离开，不依赖人工重配置。 & 需要 bootstrap、邻居发现和健康检查。 \\
可扩展 & 节点数增加时，单节点状态和每次查询成本不应线性增长。 & 通过局部路由表和分层索引实现。 \\
容错 & 节点或链路失败后仍能定位副本。 & 需要冗余路径、复制和故障检测。 \\
安全/信任 & 匿名、认证、滥用控制与隐私不能被“去中心化”自动解决。 & 需要身份、信誉和速率限制。 \\
\bottomrule
\end{tabularx}
\end{frame}

\begin{frame}{Napster：中心目录 + 端到端传输}
\begin{itemize}
  \item 中央服务器维护“用户地址—共享文件”目录；
  \item 客户端向目录查询，得到持有文件的其他用户地址；
  \item 实际文件传输由两个 peer 直接完成；
  \item 目录简化搜索，却成为单点故障、审查和规模瓶颈。
\end{itemize}
\begin{center}
\begin{tikzpicture}[>=Latex,node distance=12mm and 18mm]
\nodebox[minimum width=2.3cm]{ADSBlue}{p1}{Peer A};
\nodebox[right=of p1,minimum width=2.6cm]{ADSOrange}{dir}{中心目录};
\nodebox[right=of dir,minimum width=2.3cm]{ADSGreen}{p2}{Peer B};
\draw[\arrow,ADSGray] (p1)--node[above,font=\scriptsize]{查询}(dir);
\draw[\arrow,ADSGray] (dir)--node[above,font=\scriptsize]{地址}(p1);
\draw[\arrow,ADSBlue] (p1)--node[below,font=\scriptsize]{文件传输}(p2);
\end{tikzpicture}
\end{center}
\end{frame}

\sectionframe{二、覆盖网络与索引}{应用层覆盖图决定“下一跳是谁”，数据索引决定“对象属于谁”。}

\begin{frame}{结构化与非结构化覆盖网络}
\begin{tabularx}{\linewidth}{>{\bfseries}p{3.2cm}XX}
\toprule
结构化 overlay & 节点与对象按严格规则映射，查找路径可分析。 & 快速 lookup；复杂关键词/范围查询困难。 \\
非结构化 overlay & 邻居关系松散，对象位置不确定。 & 支持语义搜索；洪泛会造成高消息开销。 \\
\bottomrule
\end{tabularx}
\vfill
\callout{选择题}{对象名是连续数值、查询频繁且需要稳定延迟时，更适合结构化 DHT；需要模糊关键词时，非结构化搜索更灵活。}
\end{frame}

\begin{frame}{三种索引方式}
\begin{enumerate}
  \item 集中式索引：查询简单，但目录是中心瓶颈；
  \item 分布式索引：键和值映射到多个 peer，典型是 DHT；
  \item 本地索引：每个 peer 只知道自己的对象，靠 flooding 或 random walk 搜索。
\end{enumerate}
\vfill
\smallnote{语义索引支持关键词、范围和近似查询；哈希索引牺牲语义，换取确定性和均衡。}
\end{frame}

\begin{frame}{简单 DHT 的共同键空间}
\begin{center}
\resizebox{.9\linewidth}{!}{%
\begin{tikzpicture}[>=Latex]
\draw[ADSGray,thick] (0,0) rectangle (11,1.0);
\foreach \x/\lab in {0/0,2/20,4/40,6/60,8/80,11/100}{\draw[ADSGray] (\x,0)--(\x,-.15) node[below,font=\scriptsize]{\lab};}
\node[font=\scriptsize] at (5.5,.5){共同的逻辑标识空间};
\node[draw=ADSBlue,fill=ADSLightBlue,rounded corners] at (2,.5) {N2};
\node[draw=ADSOrange,fill=ADSOrange!15,rounded corners] at (6,.5) {N6};
\node[draw=ADSGreen,fill=ADSGreen!15,rounded corners] at (9,.5) {N9};
\node[font=\scriptsize,anchor=north] at (2,-.45){节点 ID};
\node[font=\scriptsize,anchor=north] at (8,-.45){对象 key};
\end{tikzpicture}}
\end{center}
\vspace{-3mm}
\begin{itemize}
  \item 节点地址与对象值都映射到同一个逻辑空间；
  \item 规则决定对象的 owner，查询不需要知道物理 IP；
  \item 节点变化时只迁移受影响的一小部分键，才有扩展性。
\end{itemize}
\end{frame}

\begin{frame}{一致性哈希的迁移直觉}
\begin{itemize}
  \item 普通哈希：节点数量变化会让大量 key 重新映射；
  \item 一致性哈希：节点和 key 都落在环上，key 归属顺时针遇到的第一个节点；
  \item 加入/离开通常只影响相邻区间；
  \item 虚拟节点可以改善负载均衡，但增加路由和故障处理状态。
\end{itemize}
\end{frame}

\sectionframe{三、Chord：环上的可扩展查找}{从顺序 successor 路由开始，再用 finger table 把 O(n) 降到 O(log n)。}

\begin{frame}{Chord 逻辑环}
\begin{center}
\begin{tikzpicture}[>=Latex,scale=.9]
\draw[ADSGray,thick] (0,0) circle (2.5cm);
\foreach \a/\lab in {20/N5,80/N18,140/N28,205/N63,260/N73,320/N99}{
  \node[draw=ADSBlue,fill=ADSLightBlue,circle,inner sep=2pt,font=\scriptsize] at (\a:2.5cm) {\lab};}
\node[font=\scriptsize] at (0,0){key 顺时针归属 successor};
\draw[\arrow,ADSOrange,thick] (20:2.5cm) arc (20:80:2.5cm);
\draw[\arrow,ADSOrange,thick] (80:2.5cm) arc (80:140:2.5cm);
\end{tikzpicture}
\end{center}
\begin{itemize}
  \item key $k$ 存在第一个节点 ID $\geq k$ 的 successor 上；
  \item 逻辑环把全局定位问题变成局部邻居关系。
\end{itemize}
\end{frame}

\begin{frame}{Chord 简单查找}
\begin{enumerate}
  \item 当前节点 i 检查 key 是否位于 $(i,successor]$；
  \item 若是，返回 successor；
  \item 否则把请求转给 successor；
  \item 最坏需要沿环走 O(n) 跳，但每个节点只需 O(1) 状态。
\end{enumerate}
\vfill
\callout{为什么先讲这个版本}{它暴露了正确性不等于性能：只要 successor 链正确，查找就能结束，但长尾会非常明显。}
\end{frame}

\begin{frame}{Finger table 的几何直觉}
\begin{center}
\begin{tikzpicture}[>=Latex,node distance=9mm]
\nodebox[minimum width=2.3cm]{ADSBlue}{n5}{节点 N5};
\node[below=of n5,align=left,font=\scriptsize] (ft) {finger[1]：$5+1$\\finger[2]：$5+2$\\finger[3]：$5+4$\\finger[4]：$5+8$\\finger[5]：$5+16$};
\nodebox[right=of ft,minimum width=2.3cm]{ADSOrange}{n18}{N18};
\nodebox[right=of n18,minimum width=2.3cm]{ADSGreen}{n63}{N63};
\draw[\arrow,ADSGray] (n5)--(n18);
\draw[\arrow,ADSGray] (n5)--(n63);
\end{tikzpicture}
\end{center}
\begin{itemize}
  \item 第 $j$ 个 finger 指向距离约 $2^{j-1}$ 的 successor；
  \item 每次选择“不超过目标的最远 finger”，跳过越来越大的区间；
  \item 路由表 O(log n)，查找平均/高概率 O(log n)。
\end{itemize}
\end{frame}

\begin{frame}{Chord 可扩展查找伪代码}
\begin{block}{LocateSuccessor(key)}
若 $key\in(i,successor]$，返回 successor；否则选择 $j=ClosestPrecedingNode(key)$，递归调用 $j.LocateSuccessor(key)$。
\end{block}
\begin{block}{ClosestPrecedingNode(key)}
从最大 finger 到最小 finger 扫描，返回落在 $(i,key)$ 内最靠后的节点。
\end{block}
\vfill
\smallnote{正确性依赖 finger 即使过期也仍能沿 successor 链前进；性能依赖周期性修复。}
\end{frame}

\begin{frame}{Chord 完整查找：N5 如何定位 key=70}
\small
标识空间为 $0\ldots127$，在线节点为 $N5,N18,N28,N63,N73,N99$。
\begin{center}
\begin{tabularx}{.94\linewidth}{>{\bfseries}p{1.4cm}p{1.8cm}X}
\toprule
步骤 & 当前节点 & 决策 \\
\midrule
1 & N5 & successor=N18，但 $70\notin(5,18]$；选择小于 70 的最远 \mbox{finger=N63}。 \\
2 & N63 & successor=N73，且 $70\in(63,73]$。 \\
3 & 返回 N73 & key=70 由第一个顺时针 ID 不小于 70 的节点 N73 负责。 \\
\bottomrule
\end{tabularx}
\end{center}
\begin{itemize}
  \item 若只沿 successor：N5$\rightarrow$N18$\rightarrow$N28$\rightarrow$N63$\rightarrow$N73，共 4 跳；
  \item 使用 finger：N5$\rightarrow$N63$\rightarrow$N73，只需 2 跳；
  \item finger 过期时仍应回退到 successor 链保证正确性。
\end{itemize}
\end{frame}

\begin{frame}{Chord 加入流程}
\begin{enumerate}
  \item 新节点 i 通过 bootstrap 找到已有节点 j；
  \item i 调用 j.LocateSuccessor(i)，得到自己的 successor；
  \item Stabilize 读取 successor 的 predecessor；
  \item Notify 让邻居更新 predecessor，并转移受影响的键；
  \item FixFingers 逐步刷新路由表。
\end{enumerate}
\end{frame}

\begin{frame}{Chord 加入的中间状态}
\begin{center}
\resizebox{.95\linewidth}{!}{%
\begin{tikzpicture}[>=Latex,node distance=18mm]
\nodebox[minimum width=1.8cm]{ADSBlue}{j}{j};
\nodebox[right=of j,minimum width=1.8cm]{ADSOrange}{i}{新节点 i};
\nodebox[right=of i,minimum width=1.8cm]{ADSGreen}{k}{k};
\draw[\arrow,ADSGray] (j)--node[above,font=\scriptsize]{successor}(i);
\draw[\arrow,ADSGray] (i)--node[above,font=\scriptsize]{successor}(k)(k);
\node[below=of i,align=center,font=\scriptsize]{Join 后：信息可能不对称\\Stabilize/Notify 后恢复环的局部不变量};
\end{tikzpicture}}
\end{center}
\end{frame}

\begin{frame}{Chord 的 churn 与故障}
\begin{itemize}
  \item Stabilize 周期性检查 successor 的 predecessor；
  \item CheckPredecessor 清理失效 predecessor；
  \item successor list 允许绕过一个或多个失败节点；
  \item finger 过期不一定破坏正确性，但会增加 hop 数和尾延迟；
  \item 复制 key 与心跳检测共同决定可用性。
\end{itemize}
\end{frame}

\begin{frame}{Chord 的复杂度与取舍}
\begin{tabularx}{\linewidth}{>{\bfseries}p{3.2cm}XX}
\toprule
本地状态 & successor、predecessor、O(log n) finger。 & 状态越多，修复成本越高。 \\
查找 & 平均/高概率 O(log n) 跳。 & 每跳都受跨地域延迟影响。 \\
加入/离开 & 只迁移邻近区间的 key。 & 短时间内会出现副本和路由不一致。 \\
负载 & 理想哈希近似均衡。 & 热点 key 仍需复制或拆分。 \\
\bottomrule
\end{tabularx}
\end{frame}

\sectionframe{四、CAN 与 Tapestry：不同的逻辑几何}{环不是唯一答案；空间划分和前缀路由也能提供可扩展查找。}

\begin{frame}{CAN：d 维 torus 空间}
\begin{center}
\begin{tikzpicture}[>=Latex]
\draw[ADSGray,thick] (0,0) rectangle (9,4.5);
\draw[ADSGray] (3,0)--(3,4.5) (6,0)--(6,4.5) (0,2.25)--(9,2.25);
\node at (1.5,3.4){N1};\node at (4.5,3.4){N2};\node at (7.5,3.4){N3};
\node at (1.5,1.1){N4};\node at (4.5,1.1){N5};\node at (7.5,1.1){N6};
\fill[ADSOrange] (6.2,1.4) circle (2pt);
\draw[\arrow,ADSBlue,thick] (1.5,1.1)--(4.5,1.1)--(7.5,1.1);
\node[font=\scriptsize] at (4.5,-.5){key 映射到坐标点，按邻居区域贪心路由};
\end{tikzpicture}
\end{center}
\vspace{-5mm}
\begin{itemize}
  \item 每个节点负责一块坐标区域；
  \item 节点只维护相邻区域，状态与维度 d 有关，而非与总节点数线性相关；
  \item 多路径和多套独立坐标空间（reality）可以提高容错。
\end{itemize}
\end{frame}

\begin{frame}{CAN 的加入与维护}
\begin{enumerate}
  \item bootstrap 返回已有节点；
  \item joiner 选择随机点并路由到该点 owner；
  \item owner 分裂自己的区域，把一半键交给 joiner；
  \item 双方更新邻居并发送 heartbeat；
  \item 故障后邻居通过 TAKEOVER 暂时接管区域，后台再做重分配。
\end{enumerate}
\end{frame}

\begin{frame}{CAN 的性能公式}
\begin{itemize}
  \item 平均邻居数约为 $O(d)$；
  \item 平均路径长度约为 $\frac d4 n^{1/d}$；
  \item 增加维度 d 可缩短路径，却增加本地状态和维护消息；
  \item 多套 reality 复制数据和邻居关系，换取可用性与多路径。
\end{itemize}
\vfill
\callout{工程判断}{若跨地域 hop 的网络延迟远大于本地 CPU，最短 hop 数不一定等于最低端到端延迟。}
\end{frame}

\begin{frame}{Tapestry：前缀逐步变长}
\begin{itemize}
  \item 节点 ID 与对象 GUID 通过哈希映射到共同空间；
  \item 沿路由每一跳增加与目标 GUID 的前缀匹配长度；
  \item 若没有同 ID 节点，选择 surrogate root；
  \item 每个路由表按 prefix level 组织，支持 O(log N) 级别查找。
\end{itemize}
\begin{center}
\begin{tikzpicture}[>=Latex,node distance=14mm]
\nodebox[minimum width=2.0cm]{ADSBlue}{t1}{6****};
\nodebox[right=of t1,minimum width=2.0cm]{ADSOrange}{t2}{62***};
\nodebox[right=of t2,minimum width=2.0cm]{ADSGreen}{t3}{62C**};
\nodebox[right=of t3,minimum width=2.0cm]{ADSGray}{t4}{62C35};
\draw[\arrow,ADSGray] (t1)--(t2)--(t3)--(t4);
\end{tikzpicture}
\end{center}
\end{frame}

\sectionframe{五、搜索、复制与 churn}{定位算法的性能，必须和对象热度、复制策略及故障模型一起看。}

\begin{frame}{非结构化搜索：flooding 与 random walk}
\begin{columns}[T,onlytextwidth]
\column{.48\linewidth}
\blocktitle{Flooding}
\begin{itemize}
  \item 查询向所有邻居扩散；
  \item TTL 或 hop count 限制范围；
  \item 成功率高，但消息重复多。
\end{itemize}
\column{.48\linewidth}
\blocktitle{Random walk}
\begin{itemize}
  \item 选择一个或多个随机邻居继续；
  \item 消息少，覆盖慢；
  \item 多 walker 往往优于单次 flooding。
\end{itemize}
\end{columns}
\end{frame}

\begin{frame}{搜索评价指标}
\begin{itemize}
  \item 成功率：找到目标对象的概率；
  \item hop/延迟：到达结果需要多少跳；
  \item node coverage：访问了多少不同节点；
  \item recall：满足关键词、范围或属性条件的对象比例；
  \item message efficiency：每条消息带来的有效召回；
  \item duplication：消息数与访问节点数之间的重复程度。
\end{itemize}
\end{frame}

\begin{frame}{对象热度与 Zipf 分布}
\begin{itemize}
  \item 查询频率常呈 Zipf-like：少数对象非常热门；
  \item 均匀复制给每个对象相同副本数，无法缓解热点；
  \item 按热度比例复制可降低热门对象查找成本，但冷门对象风险更高；
  \item 平方根复制在热点和容量之间折中。
\end{itemize}
\vfill
\callout{AI 映射}{热门模型、tokenizer 和 KV cache 的复制策略，也是在“查询热度—存储容量—跨域流量”之间做同一个优化。}
\end{frame}

\begin{frame}{复制策略比较}
\begin{tabularx}{\linewidth}{>{\bfseries}p{3cm}XX}
\toprule
Uniform & $r_i=R/m$。公平、简单；热点利用率低。 & 对象热度未知时的基线。 \\
Proportional & $r_i\propto q_i$。热门对象复制多。 & 热度估计错误时容易浪费容量。 \\
Square-root & $r_i\propto\sqrt{q_i}$。在延迟与容量间折中。 & 适合查询频率长尾。 \\
\bottomrule
\end{tabularx}
\end{frame}

\begin{frame}{churn：加入、离开、失效不是一回事}
\begin{itemize}
  \item 正常离开：可主动转移 key、副本和邻居关系；
  \item 崩溃：没有最后一次握手，邻居要靠 heartbeat/timeout 判断；
  \item 网络分区：节点可能仍在工作，却被其他节点误判为失败；
  \item 快速 churn：修复尚未完成，下一次变更又破坏旧状态。
\end{itemize}
\end{frame}

\sectionframe{六、自稳定：从任意状态恢复}{系统不要求初始状态正确，而要求最终进入合法状态并保持。}

\begin{frame}{为什么需要 self-stabilization？}
\begin{itemize}
  \item 消息丢失、内存损坏、节点恢复可能留下非法全局状态；
  \item 管理员无法穷举每种故障后的修复顺序；
  \item 自稳定算法从任意状态出发，在有限时间内收敛；
  \item 进入合法状态后，后续每一步仍保持合法。
\end{itemize}
\vfill
\callout{直觉}{不是保证“永远不坏”，而是保证“坏了以后能自己回来”。}
\end{frame}

\begin{frame}{状态机模型}
\begin{center}
\begin{tikzpicture}[>=Latex,node distance=13mm and 16mm]
\nodebox[minimum width=2.2cm]{ADSBlue}{p1s}{P1 状态};
\nodebox[right=of p1s,minimum width=2.2cm]{ADSOrange}{q1s}{P2 状态};
\nodebox[below=of p1s,minimum width=2.2cm]{ADSGreen}{q12}{消息队列 Q12};
\nodebox[below=of q1s,minimum width=2.2cm]{ADSGray}{q21}{消息队列 Q21};
\draw[\arrow,ADSGray] (p1s)--(q12);
\draw[\arrow,ADSGray] (q1s)--(q21);
\draw[\arrow,ADSGray] (q12)--(q1s);
\draw[\arrow,ADSGray] (q21)--(p1s);
\end{tikzpicture}
\end{center}
\begin{itemize}
  \item 全局 configuration = 所有节点状态 + 消息队列；
  \item privilege 表示某节点满足局部谓词，可以执行一步 move；
  \item 调度器/daemon 决定哪个 privilege 先执行。
\end{itemize}
\end{frame}

\begin{frame}{合法状态、closure、convergence}
\begin{enumerate}
  \item 用谓词 P 定义 legitimate states；
  \item Convergence：从任意非法状态出发，有限步到达 P；
  \item Closure：一旦满足 P，任何允许的下一步仍满足 P；
  \item 对互斥而言，还要讨论无死锁、无饥饿和最终唯一 token。
\end{enumerate}
\begin{center}
\begin{tikzpicture}[>=Latex]
\draw[ADSOrange,thick,->] (0,0) ellipse (2.7cm and 1.0cm) node {任意状态};
\draw[ADSBlue,thick,->] (5,0) ellipse (2.5cm and 1.0cm) node {合法状态集合 P};
\draw[\arrow,ADSGreen,thick] (2.7,0)--node[above,font=\scriptsize]{有限时间收敛}(2.5,0);
\end{tikzpicture}
\end{center}
\end{frame}

\begin{frame}{Dijkstra 自稳定 token ring}
\begin{itemize}
  \item 环上每个进程拥有一个状态值；
  \item 合法状态：系统中恰好一个进程拥有 privilege/token；
  \item 特殊节点发现与左邻居相等时递增状态；
  \item 其他节点发现与左邻居不等时复制左邻居状态；
  \item 状态空间至少需要足够大，才能打破任意初始配置的对称。
\end{itemize}
\end{frame}

\begin{frame}{Token ring 的收敛直觉}
\begin{center}
\begin{tikzpicture}[>=Latex,scale=.88]
\draw[ADSGray,thick] (0,0) circle (2.2cm);
\foreach \a/\lab in {0/P0,60/P1,120/P2,180/P3,240/P4,300/P5}{\node[draw=ADSBlue,fill=ADSLightBlue,circle,inner sep=2pt,font=\scriptsize] at (\a:2.2cm){\lab};}
\draw[ADSOrange,very thick,->] (0:2.2cm) arc (0:60:2.2cm);
\node[font=\scriptsize] at (0,-3.0){多个局部 privilege 逐步合并为一个可循环 token};
\end{tikzpicture}
\end{center}
\begin{itemize}
  \item 不要求每一步都“看起来更接近”，而是证明坏配置不可能无限停留；
  \item 证明常依赖势函数、最坏 token 传播或不可逆状态变化。
\end{itemize}
\end{frame}

\begin{frame}{Dijkstra Token Ring 完整例子：从任意状态收敛}
\small
采用 central daemon，$P_0$ 为特殊节点；状态空间 $K\geq4$。规则：$P_0$ 在 $x_0=x_3$ 时递增；其他 $P_i$ 在 $x_i\ne x_{i-1}$ 时复制左邻居。
\begin{center}
\begin{tabular}{c|c|l}
\toprule
步骤 & 状态 $(x_0,x_1,x_2,x_3)$ & 被调度节点 \\
\midrule
0 & $(2,0,2,1)$ & 任意初始配置 \\
1 & $(2,2,2,1)$ & $P_1$ 复制 $x_0$ \\
2 & $(2,2,2,2)$ & $P_3$ 复制 $x_2$ \\
3 & $(3,2,2,2)$ & $P_0$ 发现首尾相等并递增 \\
4 & $(3,3,2,2)$ & $P_1$ 复制 $x_0$ \\
5 & $(3,3,3,2)$ & $P_2$ 复制 $x_1$ \\
6 & $(3,3,3,3)$ & $P_3$ 复制 $x_2$ \\
\bottomrule
\end{tabular}
\end{center}
\smallnote{步骤 3 以后恰好一个节点 enabled，privilege 沿环移动；closure 保证系统不会再产生多个 token。}
\end{frame}

\begin{frame}{Central daemon 与 distributed daemon}
\begin{tabularx}{\linewidth}{>{\bfseries}p{3cm}XX}
\toprule
Central daemon & 每一步只选择一个可执行节点，分析更简单。 & 现实中可能过于理想化。 \\
Distributed daemon & 多个节点可能同时执行，符合异步系统。 & 并发 moves 可能互相干扰，需要更强证明。 \\
公平性 & 每个持续拥有 privilege 的节点最终都会执行。 & 没有公平性会出现 starvation。 \\
\bottomrule
\end{tabularx}
\end{frame}

\begin{frame}{自稳定互斥：从多个 token 到一个 token}
\begin{itemize}
  \item 节点失败/恢复可能制造多个 token；
  \item 算法允许多个 privilege 暂时存在；
  \item 在有限步内，重复 token 会被合并或消失；
  \item 稳态下只有一个 token，且每个想进入临界区的节点最终都能得到它。
\end{itemize}
\vfill
\callout{连接第十讲}{分布式互斥关注请求顺序和公平；自稳定互斥额外保证任意初始状态最终恢复。}
\end{frame}

\begin{frame}{自稳定 BFS 树：Dolev--Israeli--Moran 思路}
\begin{itemize}
  \item 根节点距离固定为 0；
  \item 其他节点周期性交换 distance，选择最小邻居作为 parent；
  \item distance 更新为 parent distance + 1；
  \item 即使初始 parent 指针和距离错误，系统也会逐步修复。
\end{itemize}
\end{frame}

\begin{frame}{无指定根的 BFS：比较 root ID}
\begin{itemize}
  \item 所有节点有全局唯一 ID；
  \item 当前认为 root 更大的节点会传播自己的 root 观点；
  \item 最终最大 ID（或约定的最小 ID）成为全局根；
  \item 每个节点再选择到根的最短 parent 路径。
\end{itemize}
\end{frame}

\begin{frame}{自稳定算法的成本}
\begin{itemize}
  \item Convergence span：从任意状态到安全状态的最大 transitions；
  \item Response span：到达某个目标状态的最大 transitions；
  \item 消息成本：恢复期间额外的广播、心跳和状态交换；
  \item 代价越低通常意味着状态空间、假设或故障容忍范围更受限。
\end{itemize}
\end{frame}

\sectionframe{七、AI 基础设施中的 P2P 与自恢复}{训练数据、模型、缓存和 GPU 资源都可以看作“对象 + owner + 副本 + 路由”。}

\begin{frame}{训练数据的内容寻址}
\begin{itemize}
  \item 数据分片用内容 hash 作为 key；
  \item DHT 或对象存储索引把 key 映射到持有副本的节点；
  \item 任务根据 locality 选择最近副本，减少跨地域读取；
  \item 训练完成后 checkpoint 以不可变 digest 发布，避免标签漂移。
\end{itemize}
\vfill
\callout{映射}{DHT 负责“在哪里”，一致性/签名负责“是不是同一个版本”。}
\end{frame}

\begin{frame}{模型与 tokenizer 的热点复制}
\begin{itemize}
  \item 热门基础模型在多个 GPU 节点上预热；
  \item 冷门 adapter 采用按需加载和较少副本；
  \item replica placement 同时优化命中率、显存、带宽和故障域；
  \item 模型版本变化时，旧副本通过 epoch/version 标记逐步淘汰。
\end{itemize}
\end{frame}

\begin{frame}{KV Cache：一个动态、热点、易失的对象}
\begin{itemize}
  \item key 可以由模型版本、租户、请求前缀和 tokenizer 版本组成；
  \item 热点前缀复制到靠近推理 worker 的节点；
  \item 节点失效时允许回源重算，但需要限制惊群；
  \item 若跨租户共享 cache，必须把身份和数据隔离纳入 key。
\end{itemize}
\end{frame}

\begin{frame}{GPU 调度的 overlay 视角}
\begin{center}
\begin{tikzpicture}[>=Latex,node distance=12mm and 15mm]
\nodebox[minimum width=2.2cm]{ADSBlue}{jobg}{作业};
\nodebox[right=of jobg,minimum width=2.2cm]{ADSOrange}{schedg}{调度器};
\nodebox[below left=of schedg,minimum width=2.2cm]{ADSGreen}{g1}{GPU 节点 1};
\nodebox[below right=of schedg,minimum width=2.2cm]{ADSGray}{g2}{GPU 节点 2};
\draw[\arrow,ADSGray] (jobg)--(schedg);
\draw[\arrow,ADSBlue] (schedg)--node[left,font=\scriptsize]{放置/迁移}(g1);
\draw[\arrow,ADSOrange] (schedg)--node[right,font=\scriptsize]{心跳/容量}(g2);
\draw[ADSGray,dashed] (g1)--node[below,font=\scriptsize]{邻居 gossip}(g2);
\end{tikzpicture}
\end{center}
\begin{itemize}
  \item 节点状态是动态的，不能依赖单次全局快照；
  \item gossip、租约和自稳定状态修复可以降低调度控制面的脆弱性。
\end{itemize}
\end{frame}

\begin{frame}{近期案例：Ray/KubeRay 的控制面与数据面}
\begin{itemize}
  \item 以 Ray/KubeRay 为例，作业提交、任务元数据和 worker 执行分属不同层次；控制面负责发现与调度，数据面负责对象和计算；
  \item worker 可能被抢占、重启或迁移，任务重试不能把“旧实例仍可写入”当作默认事实；需要 epoch、幂等结果和外部持久化 checkpoint；
  \item 大对象不宜在控制面日志中复制，通常记录对象引用、版本和不可变存储位置；
  \item 这把本讲的 DHT/owner/replica 路由、自稳定恢复与前面故障检测、检查点直接连在一起。
\end{itemize}
\vfill
\callout{综合题}{一个 worker 恢复后仍持有旧对象引用：如何用版本、租约和 fencing 防止旧结果覆盖新结果？}
\end{frame}

\begin{frame}{自稳定配置：从“错误参数”回到可服务状态}
\begin{itemize}
  \item 非法状态例子：一部分节点使用旧模型版本，另一部分使用新版本；
  \item 通过 epoch、版本向量和收敛规则，节点逐步采用唯一合法配置；
  \item closure 保证配置稳定后，正常心跳不会再次制造分叉；
  \item 如果故障持续发生，要区分“可收敛”与“永远被扰动”。
\end{itemize}
\end{frame}

\begin{frame}{AI 基础设施的系统级取舍}
\begin{tabularx}{\linewidth}{>{\bfseries}p{3.2cm}XX}
\toprule
一致性 & 模型/配置需要强一致，缓存和统计可弱一致。 & 先写清“旧值会不会造成错误”。 \\
可用性 & 训练可回退/重算，在线推理更重视低尾延迟。 & 不同工作负载应使用不同副本策略。 \\
安全 & 模型、数据、租户和 GPU 是不同信任域。 & 身份、签名和最小权限贯穿全链路。 \\
恢复 & 控制面需要从异常状态自动回到稳定状态。 & 自稳定思路补充故障检测和共识。 \\
\bottomrule
\end{tabularx}
\end{frame}

\sectionframe{八、课堂综合演练}{把 17 周学到的模型、算法、故障和工程约束放到同一个系统里。}

\begin{frame}{综合题：设计一个模型分发 overlay}
\begin{block}{场景}
100 个 GPU 节点、3 个地域、模型版本频繁发布；节点会抢占、重启，热门模型请求呈 Zipf 分布。
\end{block}
\begin{enumerate}
  \item 选择 Chord、CAN 或分层目录，并说明原因；
  \item 设计 key、replica、心跳和故障转移；
  \item 说明新版本发布期间如何避免读到混合权重。
\end{enumerate}
\end{frame}

\begin{frame}{综合题：故障时间线}
\begin{enumerate}
  \item t0：模型 v7 在 3 个节点上有副本；
  \item t1：节点 A 网络分区，B 继续接收请求；
  \item t2：控制面把 A 标为 suspect，并启动副本重建；
  \item t3：A 恢复但仍持有旧版本；
  \item 讨论：故障检测、epoch、读修复和自稳定规则分别做什么？
\end{enumerate}
\end{frame}

\begin{frame}{综合题：证明 closure}
\begin{itemize}
  \item 给定合法状态：每个 key 有唯一 owner，副本版本相同；
  \item 设计一次正常 join 或 heartbeat 事件；
  \item 说明该事件之后为什么仍满足“唯一 owner”或如何标记过渡状态；
  \item 如果不能证明 closure，说明系统需要的额外 epoch/锁/日志。
\end{itemize}
\end{frame}

\begin{frame}{最终系统检查清单}
\begin{enumerate}
  \item 模型：消息是否异步？节点是否可崩溃？网络是否分区？
  \item 状态：哪些是本地真相，哪些是缓存或猜测？
  \item 路由：查找正确性由什么不变量保证？
  \item 复制：副本何时产生、失效、修复、淘汰？
  \item 安全：谁可以加入、读取、发布和删除？
  \item 恢复：任意非法状态是否最终收敛？稳定后是否保持？
\end{enumerate}
\end{frame}

\sectionframe{九、课程总复习}{分布式系统不是一组 API，而是一组可验证的假设、状态和不变量。}

\begin{frame}{17 周知识地图}
\begin{center}
\resizebox{.96\linewidth}{!}{%
\begin{tikzpicture}[>=Latex,node distance=10mm and 15mm]
\nodebox[minimum width=2.4cm]{ADSBlue}{m}{模型与时间};
\nodebox[right=of m,minimum width=2.4cm]{ADSOrange}{snap}{快照与检测};
\nodebox[right=of snap,minimum width=2.4cm]{ADSGreen}{agree}{共识与一致性};
\nodebox[below=of m,minimum width=2.4cm]{ADSGray}{recover}{检查点与恢复};
\nodebox[below=of snap,minimum width=2.4cm]{ADSGray}{share}{DSM 与缓存};
\nodebox[below=of agree,minimum width=2.4cm]{ADSBlue}{secure}{认证与安全};
\nodebox[below=of secure,minimum width=2.4cm]{ADSOrange}{p2p}{P2P 与自稳定};
\draw[\arrow,ADSGray] (m)--(snap)--(agree);
\draw[\arrow,ADSGray] (m)--(recover);
\draw[\arrow,ADSGray] (snap)--(share);
\draw[\arrow,ADSGray] (agree)--(secure);
\draw[\arrow,ADSGray] (share)--(p2p);
\end{tikzpicture}}
\end{center}
\vspace{-4mm}
\callout{统一视角}{每一章都在回答：系统状态如何表示？事件如何推进？故障如何被观察？正确性如何证明？}
\end{frame}

\begin{frame}{期末项目建议}
\begin{enumerate}
  \item 选择一个真实系统：模型仓、KV cache、GPU 调度、分布式锁或数据目录；
  \item 明确系统模型、故障模型和安全边界；
  \item 给出至少一个状态不变量与一个可复现实验；
  \item 比较两种算法/一致性策略的延迟、消息和恢复成本；
  \item 提交架构图、协议时间线、故障注入脚本和 5 分钟讲解视频。
\end{enumerate}
\end{frame}

\begin{frame}{结语：从算法到系统}
\begin{itemize}
  \item 逻辑时钟告诉我们如何描述因果；
  \item 快照、检查点和故障检测告诉我们如何观察与恢复；
  \item 共识、互斥、DSM 和认证告诉我们如何协调共享状态；
  \item P2P 与自稳定告诉我们如何在动态、部分失真的环境里继续工作；
  \item AI 基础设施只是把规模、热度、租户和故障同时放大。
\end{itemize}
\vfill
\callout{最后一句}{先写假设，再写协议；先找不变量，再谈性能；先定义失败，再讨论恢复。}
\end{frame}

\begin{frame}{综合设计的最后一道边界检查}
\small
\begin{enumerate}
  \item Chord 的 finger 过期时可以回退到 successor 链保证正确性，但 churn 期间的副本迁移、并发 join/leave 和故障恢复仍需额外协议。
  \item Dijkstra 自稳定算法的经典结论依赖环拓扑、合法 token 规则和 $K\ge n$ 等条件；换成任意网络或多个 token 后，必须重新分析 closure 与 convergence。
  \item CAN、Tapestry、Chord 的“邻居”是逻辑 overlay 邻居，不一定是物理网络近邻；AI 数据和 KV Cache 系统还要同时优化拓扑距离、带宽、容量和租户隔离。
  \item “去中心化”只改变控制位置，不会自动解决认证、滥用、数据一致性或最终可用性问题。
\end{enumerate}
\vfill
\callout{期末项目要求}{任何系统图都必须同时标出逻辑拓扑、物理故障域、状态副本和恢复证据。}
\end{frame}

\end{document}
